\section{Experimental Setup}

\subsection{Datasets}

We use three publicly available single-cell RNA-seq datasets for non-spatial evaluations~\citep{luecken2020benchmarking}:
\textbf{Pancreas} (16,382 cells, 9 technologies, 14 cell types),
\textbf{PBMC (immune)} (30k cells, 5 batches, 9 cell types),
and a \textbf{Lung} dataset (30k cells, 16 batches, 17 cell types).
For spatial evaluations, we additionally use a \textbf{non-small cell lung cancer (NSCLC)} spatial transcriptomics dataset~\citep{pentimalli2025combining} (28,804 cells, 18 cell types, 9 niches).
The original dataset comprises 6 slides; we select slide 28, subsample 30,000 cells stratified by cell type and niche label, and remove cells assigned to excluded niche labels, yielding the final 28,804 cells.
Detailed dataset statistics are provided in Appendix~\ref{app:datasets}.

\subsection{Affinity Construction}

We demonstrate scProto's flexibility with two biologically complementary affinity graphs, each encoding a distinct signal that manifests in gene expression.
For transcriptomic experiments, we use an adaptive Gaussian kernel on PCA embeddings~\citep{seacells}, which scales each cell's bandwidth to its local neighborhood distance.
This corrects for the density variability inherent in single-cell data: rare cell types occupy low-density regions where a fixed bandwidth would underestimate their local similarities, and the adaptive kernel ensures they form well-connected communities in the graph.
Combined with scProto's cross-batch alignment, which pools cells from multiple batches, this enables stable metacells to form for rare cell types that are too sparsely represented within any individual batch---a limitation of within-batch methods even when using the same affinity.
Rare cell representation is assessed via coverage and rare purity.
For spatial experiments, we apply the same kernel to spatially aggregated features---the PCA embeddings of each cell averaged with those of its spatial neighbors within a fixed radius---incorporating microenvironmental context into the affinity (see Appendix~\ref{app:spatial_affinity}).
This encodes tissue organization rather than cell-type identity, a signal that transcriptomic similarity alone cannot recover; we assess it via niche purity.

\subsection{Baselines}

We compare scProto against four baselines:
\textbf{SEACells}~\citep{seacells}, a state-of-the-art metacell method based on archetypal analysis;
\textbf{MetaQ}~\citep{metaq}, which learns discrete metacell representations end-to-end via vector quantization but lacks explicit batch correction;
\textbf{scPoli (cVAE)}, scPoli's cVAE trained without its supervised prototype loss, providing a reconstruction-only baseline that isolates the contribution of scProto's additional objectives;
and \textbf{Parametric UMAP}, which preserves data geometry by aligning pairwise cell--cell relationships in the embedding space.
For spatial experiments, we additionally evaluate \textbf{SEACells (spatial)}, trained on the same spatial affinity graph as scProto, to ensure a fair comparison.
All baselines are configured to produce the same number of metacells $K$ as scProto.
Unlike scProto, neither scPoli (cVAE) nor Parametric UMAP includes a prototype or clustering mechanism; metacell assignments for these baselines are obtained by applying $K$-means clustering to their learned embeddings.

\subsection{Evaluation Protocol}

We evaluate scProto along three axes: community structure preservation, cross-batch integration, and representation of rare cell populations.
For spatial experiments, we additionally assess capture of microenvironmental structure.

To assess whether learned metacells respect the community structure of the input affinity graph, we report \textbf{modularity} of the metacell assignments with respect to the graph.
\textbf{Batch entropy} measures the diversity of batch labels within each metacell; higher values indicate more uniform mixing of cells across batches.
\textbf{Cell-type purity} measures the homogeneity of each metacell with respect to annotated cell types.

For rare cell types (defined as those whose mean per-batch frequency falls more than one standard deviation below the mean across all cell types; see Appendix~\ref{app:metrics}), we additionally report \textbf{coverage}---whether each rare cell type is represented by at least one metacell---and the \textbf{purity} of the corresponding metacells.

For spatial experiments, we evaluate \textbf{niche purity} alongside cell-type purity, measuring whether metacells reflect local tissue organization.
Detailed metric definitions are provided in Appendix~\ref{app:metrics}.
